\section{n2. Every conductor identification between centres obstructs continuity}
\textit{Supports: ENTIRE qualified G331; own n1.}

This theorem is the new bridge. Let $R$ be any positive-dimensional
excellent REDUCED local ring, in arbitrary characteristic and with arbitrary
residue field. Let $S$ be its finite normalization, retaining EVERY local
normal centre $S_i$ above $\mathfrak m$. Set
$Y=\coprod_iY_i$, $Y_i=V_{\mathfrak aS_i}(S_i)$, and let
$\rho:Y\to X$ be restriction. Assume a full fibre contains points
$w_i\in Y_i$ and $w_j\in Y_j$ for DISTINCT centres $i\ne j$.
Then there is a signed PRIMARY model obstacle $f$ on $X$ for which
$p_R(f)$ is not upper semicontinuous, neither greatest $C_R$ nor greatest
$D_R$ minorant exists, and every canonical degree has a uniform positive
gap. No normalization source is assumed to have continuous envelopes or a
finite regular cover. The normal centres themselves can be singular and
nonisolated, and every nonminimal support is retained.

Let $\mathfrak c=(R:S)$ and choose a regular element $c\in\mathfrak c$,
as granted by G331. Each $S_i$ is an excellent normal local DOMAIN.
G331 n6 gives support-zero divisorial density on its full section, with
all Rees-DVR/dimension/Izumi/historical qualifications retained. Choose a
normalized centred divisorial $\nu\in Y_i$. Then $d=\nu(c)<\infty$.
Choose an integer $N>d$, $N\ge1$.

There is $e\in S$ that is a unit at centre $i$ and belongs to
$(\mathfrak aS_k)^N$ at EVERY other centre $k$. Indeed the contractions to
$S$ of those primary ideals have distinct maximal radicals. They and the
maximal ideal at $i$ are pairwise comaximal; Chinese remainder gives
$e\equiv1$ at $i$ and $e\equiv0$ modulo each other contracted primary
ideal. This argument also works when distinct centres lie on one global
normalization component; no global idempotent separating such centres is
assumed. Since $cS\subset R$, $ce\in R$. Define
\[
       I=(ce)+\mathfrak a^N,\qquad f(v)=v(I)-N.           \tag{4}
\]
$I$ contains a primary power and is proper because $\nu(ce)=d>0$ in
this cross-centre situation. More explicitly, $c$ cannot be a unit in $R$
when normalization has an identified fibre; hence $c\in\mathfrak m$
and every normalized centred lift gives $d>0$. Thus $I$ is actual primary.
All evaluations are finite continuous, so $f$ is a signed difference of
primary profiles and $-N\le f\le0$. Put $F=\rho^*f$. On each $Y_k$,
$k\ne i$, nonnegative coefficient values and $w(e)\ge N$ imply
$w(I)=N$, including $w(ce)=\infty$. Hence $F=0$ there. At $\nu$,
$F(\nu)=d-N=-b$ with $b=N-d>0$.

The quantitative normal inequality inside G331 n6, not merely its zero
dichotomy, states that for this fixed normalized divisorial $\nu$ there
is a finite $A\ge1$ such that
\[
                  u(\nu)\ge A u(w)\quad
       (u\in C_{S_i},\ w\in Y_i).                       \tag{5}
\]
Its application is exact: $S_i$ is excellent NORMAL local domain, hence
analytically irreducible by the grant; $\mathfrak aS_i$ is primary; a
power of its maximal ideal lies in it. The granted nonunit divisorial
order bound gives $\nu(J)\le A w(J)$ for every primary ideal $J$, units
are separately zero, and uniform comparison gives (5). Enlarging $A$ to
at least one preserves (5) because $u\le0$. Taking actual decreasing
limits extends (5) to $D_{S_i}$, with extended arithmetic $A(-\infty)
=-\infty$, without subtracting infinite quantities.

For any target $C_R$ or $D_R$ minorant $g\le f$, pullback to $S_i$ is in
the corresponding granted class. Equations (4)--(5) give
\[
        (\rho^*g)(w_i)\le-b/A=-\kappa,\qquad
        p_R(f)(v_*)\le-\kappa<0,\quad v_*=\rho(w_i)=\rho(w_j).
                                                               \tag{6}
\]
This includes bottom minorants. The value $v_*(\mathfrak c)=\infty$,
because G331 says every finite-conductor fibre is a singleton.

We prove the opposite limiting behaviour on $Y_j$ by ACTUAL primary
models downstairs. By granted semilocal primary/unit gluing choose a
boundary ideal $J\subset S$ whose local ideal at $i$ is
$(\mathfrak aS_i)^N$ and whose local ideals at all other centres are units.
Thus $q(w)=-w(J)$ is $-N$ on $Y_i$ and zero elsewhere, and $q\le F$.
For $r\ge1$ set the actual ideal of $R$
\[
                  B_r=\mathfrak cJ^r+\mathfrak a^{Nr}.    \tag{7}
\]
The first summand is contained in $R$ since the conductor is an $S$-ideal
contained in $R$. Noetherianity makes it finitely generated as an $R$-ideal.
Its extension to $S$ is exactly the displayed sum. Finite ideal generators
and their nonnegative coefficients consequently give on EVERY support
\[
  \rho^*(-v(B_r)/r)(w)
          =\max\{q(w)-w(\mathfrak c)/r,-N\}.              \tag{8}
\]
At infinite conductor value the first term is $-\infty$ and the second
survives. There is no infinity-minus-infinity operation. On $Y_i$ (8)
is $-N\le F$; on every other centre it is nonpositive and hence below
$F=0$. Surjectivity makes $-v(B_r)/r$ a target minorant. It is proper
because its value at $\rho\nu$ is $Nr>0$, and contains
$\mathfrak a^{Nr}$, so is primary. Its negative effective ample divisor
$-E_{B_r}/r$ on the ACTUAL primary blowup has positive degree on EVERY
contracted proper integral curve, before any limit.

At any $w\in Y_j$ with finite conductor value, (8) tends to zero.
Nonpositivity therefore gives
\[
             p_R(f)(\rho w)=0
             \quad(w\in Y_j,\ w(\mathfrak c)<\infty).    \tag{9}
\]
This finite-conductor locus is dense in $Y_j$: the regular $c$ has finite
value at every support-zero divisorial point, and G331's full-section
divisorial density applies to the actual domain $S_j$. Choose a net of
such points $w_\lambda\to w_j$ and put $v_\lambda=\rho w_\lambda$.
Then $v_\lambda\to v_*$, while (9) is identically zero and (6) is at
most $-\kappa$. This proves a non-USC jump of at least $\kappa>0$.
By own n1, neither greatest $C_R$ nor greatest $D_R$ minorant exists.
For every actual $D_R$ minorant $g$, upper semicontinuity gives
$\limsup g(v_\lambda)\le g(v_*)\le-\kappa$, hence
\[
  \sup_X(p_R(f)-g)\ge\kappa,\qquad
  \|p_R(f)-h_m\|_\infty\ge\kappa\quad\text{for EVERY }m\ge1. \tag{10}
\]
Differences from bottom are $+\infty$. In particular canonical primary
profiles cannot converge uniformly. The proof constructs an obstacle for
every such identified fibre rather than assuming a source-envelope gap.

By own n1's raw Lipschitz bound, any continuous signed perturbation of
size $\epsilon<\kappa/2$ still has a non-USC jump and every-minorant gap
at least $\kappa-2\epsilon$: its value at $v_*$ is at most
$-\kappa+\epsilon$, whereas its values along the net are at least
$-\epsilon$. This stability requires no source continuity.

Upper semicontinuous regularization does not repair positivity here.
Define $p^*(x)=\inf_{U\ni x}\sup_{z\in U}p_R(f)(z)$. Since $f$ is
continuous and $p_R(f)\le f$, this finite USC majorant is still below
$f$. The approaching zero-valued net gives $p^*(v_*)=0$.
Equation (5) forces $p_R(f)\le-\kappa$ on the ENTIRE image of $Y_i$.
The point $x_0=\rho\nu$ has finite conductor value and belongs to no
other centre image. The other finitely many centre images are closed, so
a neighbourhood of $x_0$ meets only the image of $Y_i$. Consequently
$p^*(x_0)\le-\kappa$. If $p^*$ were in $D_R$, its pullback to $Y_i$
would satisfy (5), but its value zero at $w_i$ and its negative value
at $\nu$ contradict that inequality. Thus $p^*\notin D_R$ (and hence
$p^*\notin C_R$). This is failure of positivity in the SAME granted
geometric class, not a choice to redefine the operator or the cone.

Artifacts: n2.tex and P12.3-conductor-centre-obstruction.tex. Rerun:
python delivery.py. The universal proof is (4)--(10), not a finite
simulation, a native verdict, or a generic-valuation substitution.
